\section{Order}
%======================================================================


\begin{comment}
    - Order relation: reflexive, antisymmetric, transitive
    - Strict order
    - upper bound, upper bounds
    - greatest element: upper bound + in the set, def + uniqueness
    - least upper bound, def + uniqueness
    - maximal element
    - monotonic functions
    - cover relation -> hasse diagram
    - dual order

    - examples:

        - Mn
        - n-crown
        - partial maps
        - natural numbers
        - rational numbers (no hasse diagram

    - closure
        - operator: def idem, monotone, extensive, example ceiling
        - closure system: subset such that for all elements there is a least
        closed element above it
        - relation of both notions
        - galois connection
\end{comment}






\subsection{Basic Definitions}
%----------------------------------------------------------------------


\begin{definition}
    \label{defPoset}
    %------------------
    The pair $(S, \le)$ of a set $S$ and a relation $\le$ is called a
    \emph{partial order}
    %-------------------
    if the relation is
    \begin{enumerate}
        \item reflexive: $x \le x$

        \item antisymmetric: $x \le y \land y \le x \implies x = y$

        \item transitive $x \le y \land y \le z \implies x \le z$
    \end{enumerate}
    for all $x, y, z \in S$.
\end{definition}







\subsection{Special Elements}
%----------------------------------------------------------------------

\begin{definition}
    \label{defMinimalMaximalElement}
    %--------------------------------------------------

    An element $m \in A \subseteq P$ is a \emph{minimal element} in a subset $A$
    of a partial order $P$ if there is no strictly smaller element in the set
    $A$ i.e. $x \le a \implies x = a$ for all $x \in A$.


    An element $m \in A \subseteq P$ is a \emph{maximal element} in a subset $A$
    of a partial order $P$ if there is no strictly greate element in the set
    $A$ i.e. $a \le x \implies x = a$ for all $x \in A$.

\end{definition}



\begin{definition}
    \label{defUpperLowerBound}
    %------------------------------------------------------------

    Let $P$ be a partially ordered set and $A \subseteq P$ an arbitrary subset
    of $P$.

    An element $u \in P$ is called an \emph{upper bound} of $A$ if all elements
    in $A$ are less equal $u$ i.e. $\forall a\in A: a \le u$.

    An element $l \in P$ is called an \emph{lower bound} of $A$ if all elements
    in $A$ are greater equal $l$ i.e. $\forall a\in A: l \le a$.

    The set $A^u$ is the set of all upper bounds of $A$ and $A^l$ is the set of
    all lower bounds of $A$ i.e.
    $$
    \vertlist{
        A^u &:=& \set{x \in P: \forall a \in A: a \le x}
        \\
        A^l &:=& \set{x \in P: \forall a \in A: x \le a}
    }
    $$
\end{definition}



\begin{definition}
    \label{defMaximumMinimum}
    %--------------------------------------------------

    An element $m \in A$ of a subset $A \subseteq P$ of a partial order $(P,
    \le)$ is called a \emph{maximum} of $A$ if it is an upper bound of $A$ i.e.
    $m \in A^u$ and it is in the set $m \in A$.

    $m$ is called a \emph{minimum} of the set $A$ if it satisfies $m \in A^l$
    and $m \in A$.
\end{definition}

A maximum is in the intersection $A \cap A^u$, a minimum is in the intersection
$A \cap A^l$. Both intersections might be empty i.e. a maximum or a minimum of a
set $A$ might not exist.

If a maximum or a minimum of a set exists, then it is unique.

\begin{theorem}
    \label{thmMaximumMinimumUnique}
    \emph{Let $A \subseteq P$ of the partial order $(P, \le)$. A maximum or a
    minimum of the set $A$ (if it exists) is unique.
    }

    \begin{proof}
        Let $m, n \in P$ be two maxima of the set $A$. We have to show that both
        are equal. Since both are maxima we have $m, n \in A$ and $m,n \in A^u$.
        By definition~\ref{defUpperLowerBound} of $A^u$ we have $a \le m$ and $a
        \le n$ for all $a \in A$. This implies $m \le n$ and $n \le m$. An order
        relation is antisymmetric~\ref{defPoset} i.e. $m = n$ is valid.

        In the same manner the uniqueness of a minimum can be shown.
    \end{proof}
\end{theorem}


\begin{definition}
    \label{defSupremumInfimum}
    %--------------------------------------------------
    Let $A \subseteq P$ of the partial order $(P, \le)$.

    An element $m \in P$ is a \emph{supremum} or \emph{join} or \emph{least
    upper bound} of the set $A$ if $m$ is the minimum of $A^u$.

    An element $m \in P$ is an \emph{infimum} or \emph{meet} or \emph{greatest
    lower bound} of the set $A$ if $m$ is the maximum of $A^l$.

    Note that neither a minimum of $A^u$ nor a maximum of $A^l$ might exist.
    Therefore a join or a meet of a set $A$ might not exist.

    However of a join or a meet exists then it is unique because the maximum and
    minimum are unique~\ref{thmMaximumMinimumUnique}.

    If a join of a set $A$ exists it is written as $\bigjoin A$. If a meet of
    $A$ exists it is written as $\bigmeet A$. If the join or meet of the two
    element set $\set{a,b}$ exists we write $a \join b$ as the join of
    $\set{a,b}$ and $a \meet b$ as the meet of $\set{a,b}$
    %
    \footnote{In many texts the join of two elements is written as $a \lor b$
    and the meet as $a \land b$. Here this notation is not used to avoid
    confusion with the corresponding logical operators.}.
\end{definition}


\paragraph{Examples} Let's discuss the above definitions by looking at the
partial order

\begin{tikzpicture}
    \graph [branch right=2cm, grow up] {
        0 [xshift=1cm]
        -- {
            c/$c$ -- a,
            d -- b,
            c -- b,
            d -- a
        }
        %-- 1 [xshift=1cm];
    };
\end{tikzpicture}

\begin{enumerate}
    \item The upper bounds and the lower bounds of the empty set is the set of
        all elements of the partial order i.e.  $\emptyset^u = \emptyset^l =
        \set{0,a,b,c,d}$.

    \item For all singleton sets $\set{x}$ the element is contained in the
        upperbounds $x \in \set{x}^u$ and in the lower bounds $x \in \set{x}^l$
        of the singleton set. Therefore alll singleton sets have a maximum and a
        minimum. By definition the element is a singleton set is a minimal and a
        maximal element of the singleton set.

    \item The set $\set{a,b}$ does not have upper bounds i.e. $\set{a,b}^u =
        \emptyset$. But lower bounds exist $\set{a,b}^l = \set{c,d,0}$. Since
        $\set{a,b}^u$ is empty, the set $\set{a,b}$ does not have a join. The
        set of lower bounds is not empty, however $\set{c,d,0}$ does not have a
        maximum. Therefore $\set{a,b}$ does not have a meet either.

    \item
        Both $a$ and $b$ are maximal elements of the set of all elements.

    \item The set $\set{c,d}$ does have lower bound. It is a singleton set
        $\set{c,d} = \set{0}$. Since the element of a singleton set is the
        maximum and the minimum of the singleton set we conclude that the meet
        of $\set{c,d}$ exists i.e. we have $\bigmeet \set{c,d} = c \meet d = 0$.

    \item The element $0$ is a minimal element of the set of all elements. It is
        a minimum of the set of all elements as well.
\end{enumerate}






\subsection{Limit Element}
%----------------------------------------------------------------------

In infinite posets it is possible that an element $c$ has infinitely many
elements below it.

\begin{tikzpicture}
    \graph [grow up, branch right=2cm] {
        a0/$a_0$ -- a1/$a_1$ -- a2/$a_2$ --[dotted] c/$c$[xshift=1cm],
        b0/$b_0$ -- b1/$b_1$ -- b2/$b_2$ --[dotted] c
    };
\end{tikzpicture}

I.e. going upwards with a finite number of steps it is not possible to reach a
limit element. There are infinitely many steps needed to reach a limit element
from below. I.e. if you have an element $a_i < c$ then there exists always
another bigger element $a_{i+1}$ such that $a_i < a_{i+1} < c$.


If we take the natural numbers $\Nat$ and add a top element to it i.e. we use the
set $\set{0,1,2, \ldots, \infty}$ and use the usual order, then $\infty$ is a
limit element.

\begin{definition}
    \label{defLimitElement}
    %--------------------------------------------------

    An element $a$ is a limit element in a partial order if
    \begin{enumerate}
        \item it is not a minimal element of the order.

        \item for all $x < a$ there exists an $y$ with $x < y < a$
            i.e. $a$ has no immediate predecessor.
    \end{enumerate}

\end{definition}


\begin{theorem}
    \label{thmLimitElementIsJoin}
    \emph{A limit element $a$ in a partial order is the join of all elements
    strictly below it $a = \bigjoin \set{x: x < a}$.}

    \begin{proof}
        Let $A := \set{x: x < a}$.
        $a$ is certainly an upper bound for all elements strictly below it i.e
        $a \in A^u$. We have to prove that $a$ is the least element in $A^u$
        i.e. for all $b \in A^u$ we have  $a \le b$.

        Suppose there is a $b \in A^u$ with $b < a$. Then by definition of a
        limit element there exists an $x \in A$ with $b < x < a$. This
        contradicts the assumption that $b$ is an upper bound for all elements
        in $A$. Therefore such a $b$ cannot exist.
    \end{proof}
\end{theorem}
